\documentclass[10pt,a4paper]{amsart}
\usepackage[margin=19mm]{geometry}
\usepackage{amsmath,amssymb,mathtools}
\usepackage[scr=rsfs,cal=euler]{mathalfa}
\usepackage{xfrac}
\usepackage[unicode,hidelinks,bookmarks=false]{hyperref}
\newcommand{\ie}{i.e.\ }
\newcommand{\cf}{cf.\ }
\newcommand{\resp}{respectively\ }
\newcommand{\Subsection}{\subsection}
\providecommand{\nobreakdash}{\nobreak\hbox{-}}
\setlength{\emergencystretch}{2em}
\multlinegap0pt
\usepackage{bm}
\usepackage{bbm}
\usepackage{dsfont}
\usepackage{xspace}
\usepackage{cancel}
\usepackage{mathtools}
\usepackage{amsthm}
\makeatletter
\@ifundefined{theorem}{\newtheorem{theorem}{Theorem}}{}
\@ifundefined{lemma}{\newtheorem{lemma}[theorem]{Lemma}}{}
\makeatother
\newcommand{\Scal}{\mathcal S}
\newcommand{\Tcal}{\mathcal T}
\title[Gap spectra and densities of slow Fibonacci walks]{Gap spectra and densities of slow Fibonacci walks}
\author{Yaping Mao, Qinghong Zhao}
\date{Source: arXiv:2608.19886v1 (CC BY 4.0)}
\begin{document}
\maketitle

\noindent\emph{Source and licence.} Gap spectra and densities of slow Fibonacci walks. Version arXiv:2608.19886v1 (CC BY 4.0), distributed under the Creative Commons Attribution 4.0 International licence, as recorded in the arXiv licence metadata for this exact version and shown on the abstract page https://arxiv.org/abs/2608.19886v1. Full rights notice: licenses/arxiv-2608.19886-CC-BY-4.0.txt.\par\smallskip

\noindent\textit{Editorial scope.} Counted unit: the Theorem whose source label is \texttt{thm:phase-density} in the article (source0.tex lines 150--158, source label \texttt{thm:phase-density}), delivered together with the article's own complete original proof of it (source0.tex lines 829--940, one \texttt{proof} environment of 112 lines). No mathematics is written, repaired or re-derived: statement and proof are the article's own text line for line. The proof is detached from the statement in the source: the article states the result at lines 150--158 and prints its proof at lines 829--940, and the proof environment's own header names the statement, so the association is the article's own. Required context retained uncounted: lem:candidate (lines 196--218); lem:layer-coding (lines 260--270); eq:phase-profile (lines 819--823). Every definition, display and label of those lines is retained unchanged. Omitted from this package and not redistributed: 1--149, 159--195, 219--259, 271--818, 824--828, 941--1174. Nothing inside a retained excerpt is omitted or altered: every retained body is delivered exactly as the article's lines are written, so no omission receipt of any kind accompanies this package. Citations: the retained excerpts carry no citation command at all, so no bibliography entry of the article is transcribed and no reference is redistributed. Reference adaptation: none was needed and none was made; every reference inside the delivered unit resolves inside it, so no unresolved reference or citation is produced. Printed slips kept literal: no printed word, symbol, hypothesis or number of the counted statement or of its proof was altered. Layout and class: the article's own class is \texttt{amsart} with packages including fontenc, lmodern, microtype, geometry, latexsym, amssymb, amsmath, graphicx, float, color, fancybox, shapepar, setspace, hyperref; the wrapper is the lane's pinned \texttt{amsart} editorial wrapper at 10pt on A4, which restates the theorem-like environments the retained text needs and adds the article's own macro definitions that the retained text needs (\texttt{\textbackslash Scal}, \texttt{\textbackslash Tcal}), with extra wrapper packages (bm, bbm, dsfont, xspace, cancel, mathtools). The delivered numbering is the unit's own. Assets: no retained span contains a figure environment, a graphics-inclusion command, a PDF-page inclusion, a table or a TikZ picture; no image file is copied into this package, the package declares no asset dependency and no image file is redistributed with it. Licence: version arXiv:2608.19886v1 of this article is distributed under the Creative Commons Attribution 4.0 International licence, as recorded in the arXiv licence metadata for this exact version and shown on the abstract page https://arxiv.org/abs/2608.19886v1. A full rights notice is delivered at licenses/arxiv-2608.19886-CC-BY-4.0.txt. Characters: every retained excerpt is pure ASCII, so the delivered engine, XeLaTeX with its default Latin Modern text font, reports no missing character.
\begin{lemma}\label{lem:candidate}
For $s,n\geq2$, define $e_s(n)=\{(-1)^s\phi n\}$ and $v_s(n)=\phi n+(-1)^{s+1}e_s(n)$.\\
$(i)$ The system 
\[
\begin{pmatrix}
F_s&F_{s-1}\\
F_{s+1}&F_s
\end{pmatrix}
\binom{a_s(n)}{b_s(n)}
=
\binom{n}{v_s(n)}
\]
has the unique integer solution $a_s(n)=\phi^{1-s}n-F_{s-1}e_s(n)$ and $b_s(n)=\phi^{-s}n+F_se_s(n)$. Moreover, $\tau(n)=s$ if and only if $1\leq a_s(n)\leq b_s(n)\leq F_s$.\\
$(ii)$ Let $q_s=F_{s-1}/F_s, \kappa_s=F_s\phi^{-s}$ and $\Tcal=\{(A,B):0\leq A\leq B\leq1\}$. For $A_s=a_s(n)/F_s$ and $B_s=b_s(n)/F_s$, define $L_s(A,B)=\bigl(A+q_sB,\ \kappa_s(\phi B-A)\bigr)$.
Then
\[
L_s(A_s,B_s)=\left(\frac{n}{F_s^2},e_s(n)\right).
\]
The maps $L_s$ converge uniformly on $\Tcal$ to the invertible map
\[
L(A,B)=\left(A+\phi^{-1}B,\frac{\phi B-A}{\sqrt5}\right).
\]
\end{lemma}

\begin{lemma}\label{lem:layer-coding}
Let $(n_t)$ be a sequence of integers with $n_t\geq2$, indexed by integers $t\to\infty$ of fixed parity. Assume that $n_t/F_t^2\to u>0$. Let $\delta_{n_t}\to z$ when $t$ is even and $1-\delta_{n_t}\to z$ when $t$ is odd. Fix $h\in\mathbb Z$ and set $p_h=\{z+h\phi^{-1}\}$, with $p_h\notin\{0,\Theta(u)\}$. For all sufficiently large
$t$,
\begin{align}
t\text{ even}:\quad&
n_t+h\in D\quad\text{if and only if}\quad p_h<\Theta(u),\label{eq:layer-code-D}\\
t\text{ odd}:\quad&
n_t-h\in U\quad\text{if and only if}\quad p_h<\Theta(u).
\label{eq:layer-code-U}
\end{align}
\end{lemma}

\begin{equation}\label{eq:phase-profile}
 \mathcal P_{\ell,m}(\lambda)
 =\frac1\lambda\int_0^\lambda g_{\ell,m}(\Theta(v))\,dv
 \qquad(\lambda>0),
\end{equation}

\begin{theorem}\label{thm:phase-density} For $X\geq1$, let $t_*(X)$ be the largest even integer $t\geq2$ such that $F_t^2\leq X$, and let $\lambda_X=X/F_{t_*(X)}^2$. For every fixed $\ell,m\geq1$ and $\phi=(1+\sqrt5)/2$, there exists a continuous, multiplicatively $\phi^4$-periodic function $\mathcal P_{\ell,m}$ such that 
\[
 \frac{|D_\ell(m)\cap[1,X]|}{X}
 =\mathcal P_{\ell,m}(\lambda_X)+o(1)\qquad{and}\qquad
 \frac{|U_\ell(m)\cap[1,X]|}{X}
 =\mathcal P_{\ell,m}(\phi^2\lambda_X)+o(1).
\]
The two normalized counting functions have $\mathcal P_{\ell,m}([1,\phi^4])$ as their common set of subsequential limits. 
\end{theorem}

\begin{proof}[\bf Proof of Theorem~\ref{thm:phase-density}]
Let $t$ tend to infinity through even integers, let
$n_t=a_tF_t+b_tF_{t-1}\in\Scal_t$, and suppose that
$(a_t/F_t,b_t/F_t)\to(x,y)\in\Tcal^\circ$. By
Lemma~\ref{lem:candidate},
$(n_t/F_t^2,\delta_{n_t})=L_t(a_t/F_t,b_t/F_t)$ tends to
$L(x,y)=(u,z)$.
Assume that $\{z+h\phi^{-1}\}\notin\{0,\Theta(u)\}$ for
$0\leq h\leq m$. By Lemma~\ref{lem:layer-coding}, for all sufficiently
large $t$, $n_t+h\in D$ if and only if $\{z+h\phi^{-1}\}<\Theta(u)$ with $0\leq h\leq m$.

Let $\Omega_{\ell,m}\subset\Tcal^\circ$ consist of the points $(x,y)$
for which $c_h(x,y)=\boldsymbol1_{[0,\Theta(u))}\bigl(\{z+h\phi^{-1}\}\bigr)$ with $0\leq h\leq m$, satisfy $c_0=c_m=1$ and $\sum_{h=0}^m c_h=\ell+1$. Define $V_{\ell,m}(s)=\operatorname{area}\{(x,y)\in\Omega_{\ell,m}:u\leq s\}$.
Fix $\eta>0$. On $\Tcal\cap\{u\geq\eta\}$, the branch boundaries of
$\Theta$ and the exceptional sets $\{(x,y):\{z+h\phi^{-1}\}\in\{0,\Theta(u)\}\}$ with $0\leq h\leq m$, lie in a finite union $\mathcal L_\eta$ of line segments. Hence $\partial\Omega_{\ell,m}\cap\{u\geq\eta\}\subseteq\mathcal L_\eta\cup\partial\Tcal$.
Since the region $u<\eta$ has area $O(\eta^2)$,
$\Omega_{\ell,m}$ is Jordan measurable.

For fixed $s$, remove $\{u<\eta\}$ and the $\delta$-neighborhoods of
$\mathcal L_\eta\cup\partial\Tcal\cup\{u=s\}$. If the classification are not uniform on the remaining compact set, there would be even
integers $t_i\to\infty$, convergent coefficient grid points, and a fixed
$h$ for which it fails. By Lemma~\ref{lem:candidate}, the corresponding
scale and phase have the required limits. By
Lemma~\ref{lem:layer-coding}, this is impossible. The removed grid
proportion has limsup $O_\eta(\delta)+O(\eta^2)$. Since
$n/F_t^2=A+q_tB$ and $q_t\to\phi^{-1}$, the removal of $u=s$ also
accounts for the cutoff $n\leq sF_t^2$. Letting first $t\to\infty$ through even integers, then
$\delta\to0$, and finally $\eta\to0$, we obtain, for every fixed
$s\geq0$,
\[
\frac1{F_t^2}\#\{n\in\Scal_t:n\leq sF_t^2,\ 
n\in D_\ell(m)\}\longrightarrow V_{\ell,m}(s).
\]
Moreover, $0\leq V_{\ell,m}(s)\leq\phi^2s^2/2$, and $V_{\ell,m}$ is
continuous because every line $u=s$ has area zero. The normalized
counts and $V_{\ell,m}$ are nondecreasing in $s$ and constant for
$s\geq2$. A finite partition of $[0,2]$ on which every increment of
$V_{\ell,m}$ is less than $\varepsilon$ shows that the convergence is
uniform in $s$. For odd $t$, one has $e_t(n)=1-\delta_n$. By
Lemma~\ref{lem:candidate} and Lemma~\ref{lem:layer-coding}, $c_h$
records membership of $n-h$ in $U$. The same argument therefore gives
uniform convergence for counts of the right endpoints of occurrences
in $U_\ell(m)$.

Let $e=t_*(X)$ and write $t=e+2j$. For fixed $H$, uniformity in $s$ and
$F_{e+2j}^2/F_e^2=\phi^{4j}+o(1)$ permit termwise passage to the limit
for $|j|\leq H$. Since $1\leq\lambda_X<\phi^4+o(1)$ and
$|\Scal_t|=F_t(F_t+1)/2$, the negative tail is
$O(\sum_{j<-H}\phi^{4j})$. For $j>H$, put
$M_t=\lfloor X/F_{t-1}\rfloor$. Then
\[
 \#(\Scal_t\cap[1,X])\leq\tfrac12M_t(M_t+1),
\]
so the quadratic terms contribute $O(\sum_{j>H}\phi^{-4j})$, while
$\sum_{j>H}M_t/X\leq\sum_{j>H}F_{t-1}^{-1}=o(1)$. Both estimates are
uniform in $\lambda_X$. Letting first $X\to\infty$ and then
$H\to\infty$, we obtain
\[
 \frac{|D_\ell(m)\cap[1,X]|}{X}
 =\frac1{\lambda_X}\sum_{j\in\mathbb Z}
 \phi^{4j}V_{\ell,m}(\lambda_X\phi^{-4j})+o(1).
\]

The determinant one change of variables above has vertical section
\[
 \mathcal J(u)=\left(\frac{\phi^{-2}u}{\sqrt5},
 \min\left\{\frac{\phi^2u}{\sqrt5},1-\frac u{\sqrt5}\right\}\right)
 \quad(0<u<\phi),
\]
and $\mathcal J(u)=\varnothing$ otherwise. After setting
$v=\phi^{4j}u$, the $j$th term integrates over
$\mathcal J(v\phi^{-4j})$. For fixed $v$, these intervals are pairwise disjoint, and their union
agrees with $(0,\Theta(v))$ up to their endpoints. Indeed, if
$\xi=\phi^{4r}v\in[\phi^{-3},\phi)$, the interval for $j=-r$ has upper
endpoint $\Theta(v)$, those for $j<-r$ are empty, the remaining
intervals meet successively at their endpoints, and their lower
endpoints tend to zero. Since $\Theta(v\phi^{-4j})=\Theta(v)$ and the
event contains $z\in I_0(\Theta(v))$,
\[
 \sum_{j\in\mathbb Z}\phi^{4j}V_{\ell,m}(\lambda\phi^{-4j})
 =\int_0^\lambda g_{\ell,m}(\Theta(v))\,dv.
\]
By \eqref{eq:phase-profile}, this proves the required asymptotic formula
for $D_\ell(m)$.

The Lipschitz bound for $g_{\ell,m}$ and continuity of $\Theta$ show
that $\mathcal P_{\ell,m}$ is continuous. Substitution $v=\phi^4w$ in
\eqref{eq:phase-profile} proves its multiplicative $\phi^4$ periodicity.
For the odd layers, use $e-1$ as the base index. The same sum and tail
estimates give $\mathcal P_{\ell,m}(X/F_{e-1}^2)+o(1)$ for the
normalized right-endpoint count. After summation, replacing right
endpoints by left endpoints changes the count by at most $m$. Since
$X/F_{e-1}^2=\phi^2\lambda_X+o(1)$,
\[
 \frac{|U_\ell(m)\cap[1,X]|}{X}
 =\mathcal P_{\ell,m}(\phi^2\lambda_X)+o(1).
\]

Finally,
\[
 1\leq\lambda_X<
 \frac{F_{t_*(X)+2}^2}{F_{t_*(X)}^2}=\phi^4+o(1).
\]
Every sequence $X_k\to\infty$ therefore has a subsequence on which
$\lambda_{X_k}$ converges in $[1,\phi^4]$. Conversely, for every
$\lambda\in[1,\phi^4)$, the sequence
$X_r=\lfloor\lambda F_{2r}^2\rfloor$ satisfies
$\lambda_{X_r}\to\lambda$. Thus the down limit set is
$\mathcal P_{\ell,m}([1,\phi^4])$. The up limit set is the same as
$\mathcal P_{\ell,m}([\phi^2,\phi^6])
=\mathcal P_{\ell,m}([1,\phi^4])$.
\end{proof}

\end{document}
